\section{Pemantauan Aplikasi di Produksi: Log, Error Tracking, dan Uptime}

Setelah deployment, \textbf{pemantauan} (monitoring) memastikan aplikasi berjalan normal dan masalah terdeteksi cepat. Tiga aspek utama: (1) \textbf{log monitoring} --- periksa error log PHP, access log server, dan slow query log MySQL secara rutin; (2) \textbf{error tracking} --- alat seperti Sentry atau Bugsnag menangkap exception secara real-time dan mengirim notifikasi; (3) \textbf{uptime monitoring} --- layanan seperti UptimeRobot atau Pingdom memeriksa ketersediaan situs secara berkala dan alerting jika down \cite{mdnDeploy}.

\begin{table}[ht]
\centering
\begin{tabular}{|P{3cm}|P{4.5cm}|P{5cm}|}
\hline
\textbf{Aspek} & \textbf{Alat} & \textbf{Fungsi} \\
\hline
Error log PHP & \code{error\_log} di php.ini & Catat error PHP ke file \\
\hline
Access log & Apache/Nginx log & Lihat request masuk \\
\hline
Slow query & MySQL slow\_query\_log & Deteksi query lambat \\
\hline
Error tracking & Sentry, Bugsnag & Tangkap exception real-time \\
\hline
Uptime & UptimeRobot, Pingdom & Cek ketersediaan situs \\
\hline
APM & New Relic, Datadog & Analisis performa mendalam \\
\hline
\end{tabular}
\caption{Alat pemantauan aplikasi di produksi}
\end{table}

